\documentclass[10pt,a4paper]{amsart}
\usepackage[margin=19mm]{geometry}
\usepackage{amsmath,amssymb,mathtools}
\usepackage[scr=rsfs,cal=euler]{mathalfa}
\usepackage{xfrac}
\usepackage[unicode,hidelinks,bookmarks=false]{hyperref}
\newcommand{\ie}{i.e.\ }
\newcommand{\cf}{cf.\ }
\newcommand{\resp}{respectively\ }
\newcommand{\Subsection}{\subsection}
\providecommand{\nobreakdash}{\nobreak\hbox{-}}
\setlength{\emergencystretch}{2em}
\multlinegap0pt
\usepackage{amssymb}
\usepackage{xcolor}
\usepackage{tikz}
\usepackage[shortlabels]{enumitem}
\newtheorem{lemma}{Lemma}
\newcommand{\C}{\mathcal{C}}
\title{The structure of group-labeled graphs\\ forbidding an immersion}
\author{Rose McCarty, Caleb McFarland, Paul Wollan}
\date{Source: arXiv:2603.08820v1 (CC BY 4.0)}
\begin{document}
\maketitle

\noindent\emph{Source and licence.} The structure of group-labeled graphs\\ forbidding an immersion. Version arXiv:2603.08820v1 (CC BY 4.0), distributed by arXiv under the Creative Commons Attribution 4.0 International licence, http://creativecommons.org/licenses/by/4.0/, as recorded in the arXiv licence metadata for this exact version and shown on the abstract page https://arxiv.org/abs/2603.08820v1. Full licence notice: \texttt{licenses/arxiv-2603.08820-CC-BY-4.0.txt}.\par\smallskip

\noindent\textit{Editorial scope.} Counted unit: the article's lemma lem:uncrossingFlower (source0.tex lines 428--431). The article's complete original proof of it is delivered with it (source0.tex lines 432--450, one \texttt{proof} environment of 19 lines). No mathematics is written, repaired or re-derived: the statement and the proof are the article's own lines, in the article's own notation, and no retained line is substituted or re-flowed.  No further source-local context is needed: every cross-reference of every retained line resolves inside the two delivered spans.  Nothing was omitted from any retained span, so the package carries no omission record.  No retained line contains a citation, so the wrapper carries no bibliography and no bibliography entry.  No retained line contains a figure environment, an image-inclusion command or drawing code, so no image is delivered and the package declares no asset dependency.  Editorial formatting only, in addition: an AMS \texttt{amsart} preamble that restates the article's own declarations and macros used by the retained lines, namely the environments \texttt{lemma} on separate counters, together with the standard packages those lines need.  The retained environments print in this document as Lemma 1 in order of appearance, whereas the article prints the same items with the numbering of its own full document.  The complete native source of the article as distributed in the arXiv e-print for this exact version is bundled unchanged as \texttt{sources/196030/source0.tex}.
\begin{lemma}
\label{lem:uncrossingFlower}
    Let $\Gamma$ be a group, let $(G, \gamma)$ be a $\Gamma$-labeled graph, and let $\varphi$ be an immersion of a graph $H$ into $G$. Suppose that $H$ has a vertex which is incident to every edge of $H$, and let $x \in V(G)$ be the corresponding branch vertex of $\varphi$. Let $r$ be a positive integer and let $\Gamma'$ be a subgroup of $\Gamma$ so that $(G, \gamma)$ contains a simple flower $\C$ for $(x, \Gamma')$ of size $r$. Then we can choose $\C$ so that all but $2r$ of the branch trails of $\varphi$ are edge-disjoint from every circuit in~$\C$.
\end{lemma}

\begin{proof}
    We choose $\C$ to minimize the number of transitions of trails in $\C$ which are not transitions of $\varphi$. (It is important that every circuit in $\C$ begins at $x$; cyclically reordering a circuit changes one of its transitions.) We claim that every branch trail of $\varphi$ which does not contain the first or last edge of a circuit in $\C$ is edge-disjoint from every circuit in $\C$. Proving this claim will complete the proof of the lemma.

    Suppose otherwise, and let $T$ be a branch trail of $\varphi$ which is a counterexample. We may assume that $T$ is oriented to start from $x$. Let $\vec{e}$ be the first oriented edge of $T$ whose underlying edge $e$ is in a circuit $C \in \C$. By possibly replacing $C$ with $C^{-1}$, we may assume that $\vec{e}$ is in $C$. Let $T_1, T_2, C_1, C_2$ be trails so that $T = T_1 T_2$, $C = C_1 C_2$, and $\vec{e}$ is the first oriented edge of both $T_2$ and $C_2$. Note that by assumption $C_1$ is nonempty. Let $\vec{f}$ denote the last oriented edge of~$C_1$. Thus $(\vec{f}, \vec{e})$ is a transition of $C$ that is not a transition of~$\varphi$.

    First suppose that $T_1$ is empty. So the tail of $\vec{e}$ is $x$. Since $\gamma(C_1)\gamma(C_2) = \gamma(C) \not\in \Gamma'$, at least one of $C_1, C_2$ is labeled outside of $\Gamma'$. Both $C_1$ and $C_2$ are circuits which have $x$ as their end and do not contain the transition $(\vec{f}, \vec{e})$. Thus replacing $C$ with the appropriate choice of $C_1, C_2$ contradicts the choice of $\C$. 
    
    %RMCOMMENT{I think the notation below is more consistent with writing trails as $T = \vec{e_1}\vec{e_1} \ldots \vec{e_n}$.}
    
    Now we may assume that $T_1$ is nonempty. Note that both $C_1 T_1^{-1}$ and $T_1 C_2$ are circuits which have $x$ as their end and do not contain the transition $(\vec{f}, \vec{e})$. Also, since $\gamma(C_1 T_1^{-1})\gamma(T_1 C_2) = \gamma(C) \not\in\Gamma'$, at least one of $C_1 T_1^{-1}$ and $T_1 C_2$ is labeled outside of $\Gamma'$. Note that every transition of $T_1 C_2$ is also a transition of $T$ or $C$, including the transition ``between'' $T_1$ and $C_2$. Thus we may assume that the other circuit $C_1 T_1^{-1}$ is the one labeled outside of $\Gamma'$. To handle this case, note that $C_2$ contains some transition which is not a transition of $\varphi$ because $T$ does not contain the last oriented edge of $C_2$. Therefore replacing $C$ with $C_1 T_1^{-1}$ contradicts our choice of~$\C$. This completes the proof of Lemma~\ref{lem:uncrossingFlower}.
    %
    % Subject to this, maximize the number of branch trails in $H$ for which the first edge starting from $x$ is the first or last edge of a circuit in $\C$. Let $T$ be a branch trail in $H$ oriented to start from $x$, and let $e \in \delta(x)$ be the first edge along that trail. We claim that either $e$ belongs to a circuit in $\C$, or $T$ is edge-disjoint from every circuit in $\C$.
    %
    % Suppose otherwise, so $e$ is not in any circuit of $\C$, and let $\vec{f} \in E(G)$ be the first oriented edge of $T$ whose underlying edge is in a circuit $C \in \C$. By possibly replacing $C$ with $C^{-1}$, we may assume $\vec{f}$ is in $C$, and consider $C$ as a trail with head and tail equal to $x$. Let $T = (T_1, T_2), C = (C_1, C_2)$ where $\vec{f}$ is the first edge in $T_2, C_2$. Note that $\gamma((C_1, T_1^{-1})) + \gamma((T_1, C_2)) = \gamma(C) \not \in \Gamma'$, so one of $\gamma((C_1, T_1^{-1})), \gamma((T_1, C_2))$ is outside the subgroup. Consider the case when $\gamma((T_1, C_2)) \not \in \Gamma'$. If $C_1$ is nonempty, then the circuit $(T_1, C_2)$ does not have the transition containing $f$ and its preceding edge in $C$, and thus $(T_1, C_2)$ has fewer transitions outside those of $H$. Therefore replacing $C$ with $(T_1, C_2)$ contradicts our choice of $\C$. If $C_1$ is empty, then replacing $C$ with $(T_1, C_2)$ does not increase the transitions outside those of $H$. Furthermore this replacement causes one more trail to begin with an edge whose first or last edge is an edge of a circuit in $\C$, contradicting our choice of $\C$.
    %
    % Now suppose that $(C_1, T_1^{-1})$ is labeled outside $\Gamma'$. Note that $C_2$ can not be a subtrail of $T$ because $C$ ends at $x$ while $T$ does not. We note that if $y_i$ is a branch vertex of $H$ and $e_1, e_2$ are both the first edge along two of the branch trails from $y_i$ to the center vertex $x$, then the transition $(\{e_1, e_2\}, y_i)$ is not a transition of $H$. Thus $(C_1, T_1^{-1})$ contains fewer transitions outside those of $H$, and replacing $C$ with $(C_1, T_1^{-1})$ contradicts our choice of $\C$.
    %
    % We now claim that if the first edge $e$ along the trail $T$ starting from $x$ belongs to a circuit $C \in \C$, then in fact it is the first or last edge of $C$. Otherwise, by possibly replacing $C$ with $C^{-1}$, we may assume that $x$ is the tail of $\vec{e}$ in $C$. Let $C = (C_1, C_2)$ where $e$ is the first edge of $C_2$, and let $f$ denote the last edge of $C_1$. Then one of $C_1, C_2$ is labeled outside of $\Gamma'$, and both do not contain the transition $(\{f, e\}, x)$. Replacing $C$ with the appropriate choice of $C_1, C_2$ contradicts the choice of $\C$. Thus at most $2r$ of the branch trails of $H$ begin with the first or last edge of a circuit in $\C$, so by the above argument the remaining branch trails are edge-disjoint from all circuits in $\C$.
\end{proof}

\end{document}
